% ============================================================
% CHAPTER 4: SYSTEM AND EXPERIMENTAL SETUP
%
% Structure (matches ToC):
%   4.1  Notation and System Model
%   4.2  Debate Protocol
%        4.2.1  Agent State and Two-Phase Rounds
%        4.2.2  Memory Window and Stopping Criteria
%   4.3  Datasets
%        4.3.1  GPQA-Diamond
%        4.3.2  HiddenBench
%   4.4  Experimental Design
%
% Notation convention (consistent with Chapter 5):
%   t  = round index (0..T);  r = repetition index;  R = #repetitions
%   N  = #agents (=4);  M = #answer options (3 or 4)
%   v_i^t = vote;  c_i^t = confidence;  s_t = macrostate (vote composition)
%   \mathcal{H} = history/memory window (avoids clash with M = option count)
% Overflow (full benchmark survey, per-knob justifications) moved to Appendix.
% ============================================================

The premise of this thesis is that a multi-agent debate is a dynamical system
whose trajectory carries measurable structure. Turning that premise into empirical
claims requires a system built as a \emph{measurement instrument}, that is not exclusively 
focusing on a high-performing pipeline. Thus, the protocol must expose a
discrete vote and a confidence score per agent and per round, with private belief
cleanly separated from public message, so that the interaction log is a faithful
record of the underlying dynamics. Furthermore, the two structural axes, communication topology and 
memory window, must be treated as \emph{controlled perturbation axes}
so that we can ask which dynamiscs they cause rather than only observe.


% ------------------------------------------------------------
% 4.1 NOTATION AND SYSTEM MODEL
% ------------------------------------------------------------
\subsection{Notation and System Model}
\label{sec:setup:notation}

We study a fixed group of $N$ language-model agents that repeatedly answer the
same multiple-choice question over a sequence of discrete rounds. All agents are
instances of a single frozen model $f_\theta$; they differ only in the private
context each one has seen. Because the same symbols recur throughout the
detection methodology of \autoref{sec:method}, we fix the notation here.
\autoref{tab:notation} summarises the symbols used across this thesis.

\begin{table}[H]
\centering
\small
\begin{tabular}{ll}
\toprule
Symbol & Meaning \\
\midrule
$N$ & number of agents ($N = 4$ throughout) \\
$i, j$ & agent indices, $i \in \{1, \ldots, N\}$ \\
$t$ & round index, $t \in \{0, 1, \ldots, T\}$ \\
$T$ & hard cap on the number of rounds \\
$r$ & repetition index (an independent re-run of a task--config cell) \\
$R$ & number of repetitions per task--config cell \\
$M$ & number of answer options a task presents, $M \in \{3, 4\}$ \\
$v_i^{t}$ & vote of agent $i$ at round $t$, one of the $M$ options \\
$c_i^{t}$ & private confidence of agent $i$ at round $t$, $c_i^t \in \{1,\ldots,10\}$ \\
$\mathcal{N}(i)$ & neighbourhood of agent $i$ under the communication topology \\
$W$ & memory window (rounds of history each agent may see) \\
$u$ & early-stopping threshold (consecutive unanimous rounds) \\
$s_t$ & system macrostate at round $t$ (the vote composition, see below) \\
\bottomrule
\end{tabular}
\caption[Notation used throughout the thesis]{Core notation used from this Section onward. Rounds are indexed by $t$ and
repetitions by $r$; the two must not be confused.}
\label{tab:notation}
\end{table}

Three units of analysis must be kept apart, since they nest inside one another
and are easy to confuse. A \emph{round} is a single synchronous time step $t$
within one debate, at which every agent (re)casts its vote; a debate runs for up
to $T$ rounds. A \emph{repetition} (indexed by $r$) is one complete debate on a
given task from round $0$ to termination; because sampling is stochastic, the
same task run twice yields different repetitions. A \emph{task--config cell} is a
single question paired with one configuration (a topology and a memory window),
and we execute $R$ repetitions of each cell. In short: many rounds make up one
repetition, and many repetitions make up one task--config cell.

With the notation established, two structural ingredients complete the system model.
The \emph{interaction structure} is a binary adjacency matrix
$\mathbf{G} \in \{0,1\}^{N \times N}$, where $G_{ij}=1$ means agent $i$ receives
agent $j$'s messages. Self-loops are excluded ($G_{ii}=0$), so no agent is fed
its own output as a neighbour signal, and the neighbourhood of agent $i$ is
$\mathcal{N}(i) = \{j \mid G_{ij}=1\}$. Any topology can be instantiated by
supplying a different matrix without changing the agent or round logic. We use
two topologies:
\begin{itemize}
  \item \textbf{Fully connected:} $G_{ij} = 1 - \delta_{ij}$, so every agent
        observes all others each round. This maximises information flow and
        serves as the baseline. With $N = 4$ agents, each round involves
        $N(N-1) = 12$ messages.
  \item \textbf{Star:} one hub agent is connected to all $N-1$ spokes, and the
        spokes are connected only to the hub. The hub is drawn uniformly at
        random per repetition, independently of agent indices, so that any
        observed hub effect is structural rather than an artefact of indexing.
        Each round involves $2(N-1) = 6$ messages, half the fully-connected cost.
        Star therefore trades per-round token cost for less information flow: the hub
        is the sole information bottleneck, so agreement typically takes more rounds
        to propagate.
\end{itemize}

The second structural ingredient is the system-level state variable.
Individual agents are exchangeable: only how many of them hold each option
matters for the system-level analyses, not which agent holds which. We therefore
summarise the group at round $t$ by its \emph{vote composition}, the count of
agents on each option,
\begin{equation}
  s_t = (n_1^{t}, n_2^{t}, \ldots, n_M^{t}),
  \qquad \sum_{k=1}^{M} n_k^{t} = N,
  \label{eq:macrostate}
\end{equation}
where $n_k^{t}$ is the number of agents voting for option $k$ at round $t$. The
number of agents $N$ and the number of options $M$ are independent: $N = 4$ throughout, with $(4,0,0,0)$, $(3,1,0,0)$, $(2,2,0,0)$,
$(2,1,1,0)$, $(1,1,1,1)$ as the possible macrostates for $M = 4$ tasks. This macrostate is the object on which the
system-level motif detectors of \autoref{sec:method} operate.


% ------------------------------------------------------------
% 4.2 DEBATE PROTOCOL
% ------------------------------------------------------------
\subsection{Debate Protocol}
\label{sec:setup:protocol}

The protocol specifies three things: the persistent state each agent carries
across rounds, the within-round message passing that produces a new state, and
the memory window controlling how much past history each agent may see. We follow
the multi-agent debate paradigm of \citet{du2024improving}, in which agents
independently answer and then revise over several rounds, but instrument it so
that the interaction log exposes exactly the quantities the motif detectors
require.

\label{sec:setup:protocol:state}%
These requirements translate into a specific per-agent state. At each round $t$, agent $i$ produces four fields:
\begin{itemize}
  \item $v_i^{t}$: the discrete \textbf{vote}, one of the task's $M$ options;
  \item $c_i^{t} \in \{1,\ldots,10\}$: a private \textbf{confidence}, never transmitted;
  \item $\text{priv}_i^{t}$: private chain-of-thought \textbf{reasoning}, never transmitted;
  \item $\text{pub}_i^{t}$: a \textbf{public message}, broadcast to neighbours.
\end{itemize}
The separation of the public message from the private reasoning and confidence is
intentional as it lets us measure whether agents say what they privately hold. 
Furthermore, the split of vote and confidence follows the continuous-opinion,
discrete-action template of \citet{martins2008coda}, and the public/private
channel separation follows the DEBATE design of \citet{chuang2025debate}.

Building on this state structure, for rounds $t \geq 1$ each agent also emits an ephemeral \emph{argument} before
committing to a belief:
\begin{equation}
  \text{argument}_i^{t} = \big(\text{defense}_i^{t},\ \text{challenge}_i^{t},\ \text{question}_i^{t}\big),
\end{equation}
consisting of a \textbf{defense} (one concrete reason the current vote is
correct), a \textbf{challenge} (a dispute of a specific peer claim, or a
concession), and a \textbf{question} (a targeted query to a named peer). The
argument is visible to neighbours at the start of the \updatephase{} phase but does not itself change
any vote. These fields structure the deliberation but are not themselves used
quantitatively, but helped constructing a system of higher debate intensity. 

Together, the committed state and the ephemeral argument define the round structure.
Each round consists of two sequential phases; within each phase all agents act in
parallel, but the \arguephase{} phase runs to completion before the \updatephase{} phase begins.
Round $0$ is asymmetric, as no communication has yet occurred. \autoref{fig:protocol} illustrates the full message flow for two agents over the
first three rounds.

\begin{itemize}
  \item \textbf{Round $0$ (initialisation, \updatephase{} phase only).} Each agent receives
        the question and options and produces an initial belief with no peer
        input:
        \[
          (v_i^{0},\ c_i^{0},\ \text{priv}_i^{0},\ \text{pub}_i^{0})
          = f_\theta\big(\text{question, options}\big).
        \]

  \item \textbf{Round $t \geq 1$, \arguephase{} phase.} Agent $i$ reads its self
        and peer history windows $\mathcal{H}^{\text{self}}_i(t,W)$ and
        $\mathcal{H}^{\text{peer}}_i(t,W)$ (defined formally below)
        and emits an argument. Votes are \emph{not} updated in this phase:
        \[
          \text{argument}_i^{t}
          = f_\theta\big(q,\ \mathcal{H}^{\text{self}}_i(t,W),\ \mathcal{H}^{\text{peer}}_i(t,W)\big).
        \]

  \item \textbf{Round $t \geq 1$, \updatephase{} phase.} Agent $i$ additionally
        reads its own and its neighbours' current-round arguments, then updates its
        full state; if it changes its vote it is instructed to cite the specific
        argument that caused the change:
        \[
          (v_i^{t},\ c_i^{t},\ \text{priv}_i^{t},\ \text{pub}_i^{t})
          = f_\theta\big(q,\ \mathcal{H}^{\text{self}}_i(t,W),\ \mathcal{H}^{\text{peer}}_i(t,W),\
          \{\text{argument}_k^{t}\}_{k \in \{i\} \cup \mathcal{N}(i)}\big).
        \]
\end{itemize}

\begin{figure}[H]
\centering
\begin{tikzpicture}[
  font=\small,
  phB/.style={draw, rounded corners=3pt, fill=blue!8, align=left,
              inner sep=6pt, text width=5.3cm},
  phA/.style={draw, rounded corners=3pt, fill=orange!12, align=left,
              inner sep=6pt, text width=5.3cm},
  parr/.style={-Stealth, thick, gray!65},
  darr/.style={-Stealth, thick, orange!80},
  >=Stealth
]

\coordinate (Lsep) at (-3.8, 0);
\coordinate (Rsep) at (11.8, 0);

% Question box
\node[draw, rounded corners=3pt, fill=gray!10, align=center,
      inner sep=6pt, text width=11.2cm] (Q) at (4.0, 2.7) {
  \textbf{Question $q$ + Options ($M$ choices)}
};

% r=0 --- Phase B only
\node[phB] (B0i1) at (0, 0) {
  \textbf{Agent $1$\,\textbullet\,\updatephase{} phase}\\[2pt]
  $v_1^{0}$: vote\\[2pt]
  {\color{gray}$c_1^{0}$: confidence \textit{(private)}}\\[2pt]
  {\color{gray}$\text{priv}_1^{0}$: reasoning \textit{(private)}}\\[2pt]
  $\text{pub}_1^{0}$: public message
};
\node[phB] (B0i2) at (8.0, 0) {
  \textbf{Agent $2$\,\textbullet\,\updatephase{} phase}\\[2pt]
  $v_2^{0}$: vote\\[2pt]
  {\color{gray}$c_2^{0}$: confidence \textit{(private)}}\\[2pt]
  {\color{gray}$\text{priv}_2^{0}$: reasoning \textit{(private)}}\\[2pt]
  $\text{pub}_2^{0}$: public message
};
\node[font=\bfseries\small, left=0.5cm of B0i1] {$t{=}0$};
\draw[-Stealth, thick, gray!45] (Q.south) -- (B0i1.north);
\draw[-Stealth, thick, gray!45] (Q.south) -- (B0i2.north);

\path (B0i1.south) +(0,-0.35) coordinate (sep01);
\draw[dashed, gray!45, thin] (Lsep |- sep01) -- (Rsep |- sep01);

% r=1 --- Phase A
\node[phA, below=0.7cm of B0i1] (Ar1i1) {
  \textbf{Agent $1$\,\textbullet\,\arguephase{} phase}\\[2pt]
  $\text{argument}_1^{1}$: \textit{defense, challenge, question}
};
\node[phA, below=0.7cm of B0i2] (Ar1i2) {
  \textbf{Agent $2$\,\textbullet\,\arguephase{} phase}\\[2pt]
  $\text{argument}_2^{1}$: \textit{defense, challenge, question}
};
\node[font=\bfseries\small, left=0.5cm of Ar1i1] {$t{=}1$};
\draw[parr] (B0i1.south) -- (Ar1i2.north);
\draw[parr] (B0i2.south) -- (Ar1i1.north);

% r=1 --- Phase B
\node[phB, below=0.5cm of Ar1i1] (Br1i1) {
  \textbf{Agent $1$\,\textbullet\,\updatephase{} phase}\\[2pt]
  $v_1^{1},\ c_1^{1},\ \text{priv}_1^{1},\ \text{pub}_1^{1}$
};
\node[phB, below=0.5cm of Ar1i2] (Br1i2) {
  \textbf{Agent $2$\,\textbullet\,\updatephase{} phase}\\[2pt]
  $v_2^{1},\ c_2^{1},\ \text{priv}_2^{1},\ \text{pub}_2^{1}$
};
\draw[darr] (Ar1i1.south) -- (Br1i2.north);
\draw[darr] (Ar1i2.south) -- (Br1i1.north);

\path (Br1i1.south) +(0,-0.35) coordinate (sep12);
\draw[dashed, gray!45, thin] (Lsep |- sep12) -- (Rsep |- sep12);

% r=2 --- Phase A
\node[phA, below=0.7cm of Br1i1] (Ar2i1) {
  \textbf{Agent $1$\,\textbullet\,\arguephase{} phase}\\[2pt]
  $\text{argument}_1^{2}$: \textit{defense, challenge, question}
};
\node[phA, below=0.7cm of Br1i2] (Ar2i2) {
  \textbf{Agent $2$\,\textbullet\,\arguephase{} phase}\\[2pt]
  $\text{argument}_2^{2}$: \textit{defense, challenge, question}
};
\node[font=\bfseries\small, left=0.5cm of Ar2i1] {$t{=}2$};
\draw[parr] (Br1i1.south) -- (Ar2i2.north);
\draw[parr] (Br1i2.south) -- (Ar2i1.north);

% r=2 --- Phase B
\node[phB, below=0.5cm of Ar2i1] (Br2i1) {
  \textbf{Agent $1$\,\textbullet\,\updatephase{} phase}\\[2pt]
  $v_1^{2},\ c_1^{2},\ \text{priv}_1^{2},\ \text{pub}_1^{2}$
};
\node[phB, below=0.5cm of Ar2i2] (Br2i2) {
  \textbf{Agent $2$\,\textbullet\,\updatephase{} phase}\\[2pt]
  $v_2^{2},\ c_2^{2},\ \text{priv}_2^{2},\ \text{pub}_2^{2}$
};
\draw[darr] (Ar2i1.south) -- (Br2i2.north);
\draw[darr] (Ar2i2.south) -- (Br2i1.north);

% Legend
\coordinate (legX) at (4.0, 0);
\path (Br2i1.south) +(0,-1.15) coordinate (legY);
\node[draw, rounded corners=3pt, fill=gray!5, align=left, inner sep=6pt]
  at (legX |- legY) {%
  {\color{gray!65}$\longrightarrow$}~$\text{pub}$: public message, sent to each
    neighbour's \arguephase{} phase the \emph{next} round\\[3pt]
  {\color{orange!80}$\longrightarrow$}~$\text{argument}$: \arguephase{} phase output, sent to
    each neighbour's \updatephase{} phase the \emph{same} round\\[3pt]
  {\color{gray}$c$, $\text{priv}$}: confidence and reasoning, produced every
    \updatephase{} phase but \textbf{never transmitted}%
};

\end{tikzpicture}
\caption[Two-phase synchronous debate protocol]{Two-phase synchronous protocol illustrated for two agents over rounds
$t{=}0,1,2$. As round~$0$ has no \arguephase{} phase, the agents read
the question directly. Each \updatephase{} phase commits the four state fields
$(v,c,\text{priv},\text{pub})$; the confidence $c$ and reasoning $\text{priv}$
stay private (grey). A public message reaches neighbours in the \emph{next}
round's \arguephase{} phase; an \arguephase{} argument reaches neighbours in the \emph{same} round's
\updatephase{} phase, which is what lets one round carry a full question--answer exchange.}
\label{fig:protocol}
\end{figure}

The two-phase structure compresses a question--answer exchange into a single
round as a one-phase synchronous protocol could not achieve this without a
two-round lag. For the dynamical analysis, only the committed \updatephase{} state is
treated as the system state at round $t$; the arguments mainly function as message memory and debate facilitator. All outputs are returned as structured JSON and no problems with malformed responses were observed. The complete prompt templates are shown in \autoref{app:prompts}.

\label{sec:setup:protocol:memory}%
One aspect of the round structure left unspecified so far is how much history each
agent may access. The memory window $W \in \mathbb{Z}_{\geq 1}$ controls how many past rounds each agent sees. At round $t$, agent $i$ has a \emph{self} window of its
own recent states and arguments, and a \emph{peer} window of its neighbours' recent
votes, public messages, and arguments:
\begin{align}
  \mathcal{H}^{\text{self}}_i(t, W)
    &= \big\{\big(v_i^{t-k},\, c_i^{t-k},\, \text{pub}_i^{t-k},\, \text{argument}_i^{t-k}\big)\big\}_{k=1}^{\min(t,\,W)}, \\
  \mathcal{H}^{\text{peer}}_i(t, W)
    &= \big\{\big(v_j^{t-k},\, \text{pub}_j^{t-k},\, \text{argument}_j^{t-k}\big)\big\}_{j \in \mathcal{N}(i),\ k=1..\min(t,\,W)}.
\end{align}
Peer entries are presented in randomised order at each call to mitigate position
bias \citep{zheng2024selectors}, and round-$0$ entries omit the argument field, which does not exist yet. A larger memory window
$W$ exposes more history but also interacts with documented long-context
degradation \citep{liu2024lostmiddle}, which is why we treat $W$ as an explicit
perturbation axis rather than fixing it implicitly. $W = 1$ is the near-Markov
baseline in which each agent sees only the previous round.

Now that the protocol and memory structure are defined, the remaining question is when to stop.
Running every repetition for the full $T$ rounds is wasteful once agents have
converged: pilot transition matrices show a self-loop probability of $0.97$--$1.00$
for the unanimous states, meaning that once consensus is reached, escape is rare.
We therefore stop a repetition as soon as $u$ consecutive rounds are
unanimously agreed. The threshold $u$ trades efficiency against reliability: too
small and a transient consensus may be cut short; too large and the savings
shrink. Analysing pilot repetitions (see \autoref{fig:time_savings} in the appendix), we find that the
harmful case, in which stopping would have prevented a later recovery to the
correct answer, stays negligible while coverage remains high, and we set $u = 3$
with a hard cap of $T = 15$. 
Terminating debate at consensus is standard practice
in the MAD literature \citep{yin2023eot, yang2025revisiting, kaesberg2025voting}.

When a repetition ends, the system's answer is the \emph{plurality} of the agents' votes
at the final round, the option held by the most agents. Because a repetition stops
only after $u = 3$ consecutive unanimous rounds, every converged repetition ends
with all $N$ agents on one option, so the plurality is unambiguous and equals that
option. A tie can therefore arise only in the rare repetitions that reach the hard
cap $T = 15$ without converging, which occurs in approximately $2\,\%$ of repetitions. Such ties are broken deterministically towards the first option in
canonical order. The resulting label is a low-confidence artefact of wherever the
unresolved run happened to land rather than a genuine group decision, and
\autoref{sec:results:lc} recommends flagging these repetitions as unresolved
instead of trusting the reported majority.

% ------------------------------------------------------------
% Commented out but retained in case it is needed later.
% ------------------------------------------------------------
% \paragraph{Early task stopping.}
% Not every question is worth its full repetition budget. For some, all agents are
% already unanimous at $t = 0$, so the outcome is fixed by the model's prior rather
% than by deliberation. Using a small pilot batch of $R_{\text{pilot}} = 5$
% repetitions, we skip the remainder of a task if the mean per-agent accuracy at
% $t = 0$ and the mean system accuracy at the final round agree at an extreme:
% \begin{equation}
%   \big(\bar{a}_0 \geq \theta_{\text{high}} \ \text{and}\ \bar{a}_T \geq \theta_{\text{high}}\big)
%   \quad\text{or}\quad
%   \big(\bar{a}_0 \leq \theta_{\text{low}} \ \text{and}\ \bar{a}_T \leq \theta_{\text{low}}\big),
% \end{equation}
% with $\theta_{\text{high}} = 0.95$ (trivially easy) and $\theta_{\text{low}} = 0.05$
% (effectively unsolvable). Requiring agreement between $t = 0$ and the final round
% prevents false positives, since a question whose accuracy moves between the two
% cannot satisfy either condition. The unanimously-correct case is safe to filter
% unconditionally (a correct consensus at $t = 0$ is stable under the dynamics);
% the unanimously-wrong case is in principle riskier, but recovery from it was
% observed in only one of twenty pilot questions.


% ------------------------------------------------------------
% 4.3 DATASETS
% ------------------------------------------------------------
\subsection{Datasets}
\label{sec:setup:datasets}

Because this work studies motif emergence rather than task-completion
performance, our benchmark selection criteria differ from those usual in the MAD
literature. A benchmark suitable for motif detection must provide (i) a
\textbf{discrete shared decision variable}, so that all agents commit to a
position in the same finite set each round; (ii) \textbf{verifiable ground truth},
so dynamics can be related to correctness; (iii) \textbf{calibrated difficulty},
since saturated tasks remove the disagreement that drives deliberation; and (iv)
\textbf{sufficient initial disagreement}, without which no motif can be
distinguished from trivial consensus.

Applying these criteria rules out most existing benchmarks. Saturated reasoning
sets such as GSM8K and MMLU leave no headroom for deliberation dynamics; native
multi-agent benchmarks are largely incompatible, being either role-specialised
or defined over spatial rather than categorical positions. A full survey of the candidates we considered and the reason each was
rejected is given in \autoref{app:benchmarks}. Two bench satisfy all four
criteria and are complementary, so we use them throughout this study.

\paragraph{GPQA-Diamond.}
\label{sec:setup:datasets:gpqa}
GPQA-Diamond \citep{rein2024gpqa} is a graduate-level science
question-answering benchmark of 198 questions in biology, physics, and
chemistry, written by domain experts and validated through a ``Google-proof''
protocol that resists retrieval-based shortcuts. It satisfies all four criteria
directly: four-option multiple choice ($M = 4$) provides the discrete decision
variable; expert validation gives unambiguous ground truth; frontier models fall
well below ceiling \citep{rein2024gpqa, becker2025orchestration}, leaving genuine
headroom for debate; and agents genuinely disagree at round one
\citep{becker2025orchestration}.
GPQA-Diamond also has direct multi-agent precedent
\citep{becker2025orchestration, kenton2024debate, he2025dimo, wang2025mars}. We
adopt it as the \textbf{primary} dataset. 
The main cavueat is that this benchmark has less cumulative MAD literature compared to MMLU or GSM8K, which is however mostly the case as its a moer recent benchmark and not because of a methodological pivot. 
A second caveat concerns comparability of the accuracy numbers: published figures
refer to strong models on the full benchmark, whereas we use
a mid-tier model on a screened task subset only.
Our results are therefore not comparable to published frontier GPQA-Diamond
numbers, and are not intended to be; the screening pass deliberately selects the
most dynamically contested tasks, where accuracy is low by design.

\paragraph{HiddenBench.}
\label{sec:setup:datasets:hiddenbench}
Moreover, we use HiddenBench \citep{li2026hiddenbench} as a structurally complementary
\textbf{secondary} dataset. It formalises the
hidden-profile paradigm from social psychology
\citep{stasser1985pooling, schulzhardt2012synergy}, in which the evidence needed for the
correct answer is split across agents, so that no single agent can solve the task
alone and success requires information pooling. Its value here is as a
control for a specific confound. Motifs seen only on GPQA-Diamond might be
artefacts of cooperative reasoning over a reliable individual signal, in which
RLHF-induced truth-seeking \citep{ouyang2022instructgpt} amplifies convergence.
Under the hidden-profile condition, by contrast, pre-discussion accuracy is very
low by construction \citep{li2026hiddenbench}, so any post-discussion improvement
must come from social information transfer rather than from individual competence.
Motifs that appear on both datasets are therefore robust to this confound,
while motifs that appear on only one are themselves informative on a dataset-level.

One structural detail of HiddenBench matters for the analysis. While we use
$N = 4$ agents throughout, the number of answer options is task-dependent,
$M \in \{3, 4\}$. The agent count $N$ and the option count $M$ are therefore
independent quantities and should not be conflated. This is why every
option-indexed quantity in \autoref{sec:method} is defined over the task's actual
$M$ options.
Extension to mixed-motive negotiation datasetes is left to future work.


% ------------------------------------------------------------
% 4.4 EXPERIMENTAL DESIGN
% ------------------------------------------------------------
\subsection{Experimental Design}
\label{sec:setup:design}

The system exposes many tunable components. We classify each as \textbf{scan} (a
primary perturbation axis), \textbf{fix} (held constant), or \textbf{defer} (out
of scope). \autoref{tab:knob-summary} summarises the classification and the
literature supporting it; the detailed per-knob justification is deferred to
\autoref{app:knobs}.

Of these components, two are varied as primary perturbation axes, chosen because they
are the canonical levers for network-motif analysis:
\begin{itemize}
  \item \textbf{Memory window} $W \in \{1, 2, 5\}$, with $W = 1$ the near-Markov
        baseline. Memory content is an established driver of debate behaviour
        \citep{du2024improving, li2024sparsetopology}.
  \item \textbf{Topology} $\in \{\textsc{fc}, \textsc{star}\}$. Fully-connected
        probes the symmetric baseline; the star probes hub-induced dominance,
        following evidence that centralised structures concentrate influence
        \citep{li2024sparsetopology, chen2023consensus, liu2025problemsolving}.
\end{itemize}
Their product gives $3 \times 2 = 6$ configurations per task. The dataset (GPQA
or HiddenBench) is classified as a scanned condition in \autoref{tab:knob-summary}
because the information condition varies between benchmarks, but it is not a
structural axis in the sense of a manipulation applied within a fixed task; the
two benchmarks are instead analysed separately. The full design
therefore spans $6 \times 2 = 12$ configuration--dataset cells.
All remaining components are held constant to isolate the effect of the two scan
axes:
\begin{itemize}
  \item \textbf{Agents} $N = 4$: the smallest value at which the two
        topologies are structurally distinct and multi-way splits can occur.
  \item \textbf{Rounds} $T = 15$ (hard cap) with early stopping $u = 3$: $T$ is
        an observation budget rather than a structural knob, since trajectories
        can be cropped post hoc but not extended.
  \item \textbf{Model} \texttt{mistral-medium-instruct}: a single mid-tier
        model to isolate topology and memory effects from capability-asymmetry
        effects of heterogeneous models \citep{wynn2025talk}.
  \item \textbf{Sampling} temperature $0.7$, reasoning disabled: diverse yet
        focused outputs, avoiding the degenerate identical-sample case
        \citep{zhu2026demystifying}.
  \item \textbf{Framing} cooperative, with the full \arguephase{}/\updatephase{} structured
        output as the measurement instrument, and synchronous execution to
        remove agent-order confounds.
\end{itemize}

Finally, selecting which tasks to run completes the design. Motifs are per-task properties, and running the full grid at high $R$ on every
question is infeasible: at roughly $25$\,s per repetition, the full GPQA-Diamond
set across six configurations would take approximately $10.3$ days ($248$\,h) at
$R = 30$ or $17.2$ days ($413$\,h) at the final $R = 50$; a two-phase design that 
employs an initial screening process reduces this to approximately $29$\,h. 
It begins with a shallow \emph{screening} pass which runs every task once at a
single baseline configuration (\textsc{fc}, $W = 1$, $R = 3$) and retains only
those tasks where debate actually occurs \emph{and} the final answer varies
across repetitions, since only those tasks have dynamics, rather than the model
prior, deciding the outcome. Tasks that are unanimous at $t = 0$, that always
converge to the same answer, or where no genuine disagreement arises, are
excluded.

\begin{table}[H]
\centering
\small
\begin{tabular}{p{3.0cm}lp{4.4cm}p{4.8cm}}
\toprule
\textbf{Knob} & \textbf{Verdict} & \textbf{Value} & \textbf{Key references} \\
\midrule
Memory $W$ & scan & $\{1, 2, 5\}$ & \citep{du2024improving, li2024sparsetopology} \\
Topology & scan & fc, star & \citep{li2024sparsetopology, chen2023consensus, liu2025problemsolving} \\
Agents $N$ & fix & $4$ & \citep{du2024improving, li2026hiddenbench} \\
Rounds $T$ & fix & $15$ (cap), $u=3$ & \citep{kaesberg2025voting, yang2025revisiting} \\
Output structure & fix & full \arguephase{}/\updatephase{} & \citep{liang2024encouraging, he2025dimo} \\
Framing & fix & cooperative & \citep{liang2024encouraging, du2024improving} \\
Synchrony & fix & parallel-synchronous & \citep{du2024improving, liu2024groupdebate} \\
Model & fix & \texttt{mistral-medium-instruct} & \citep{liu2025stopovervaluing, wynn2025talk} \\
Temperature & fix & $0.7$, no reasoning & \citep{zhu2026demystifying} \\
Information cond. & scan (dataset) & GPQA ($M{=}4$), HiddenBench ($M{\in}\{3,4\}$) & \citep{li2026hiddenbench, stasser1985pooling} \\
Tool access & defer & none & \citep{rein2024gpqa} \\
\bottomrule
\end{tabular}
\caption[Classification of tunable components]{Classification of tunable components. Full justification per knob is
given in \autoref{app:knobs}.}
\label{tab:knob-summary}
\end{table}

Retained tasks are promoted to a deep \emph{characterisation} pass on the full
six-configuration grid. This yields the final dataset used throughout
\autoref{sec:results}: \textbf{50 tasks for each of the two benchmarks, six
configurations each, $R = 50$ repetitions}, for a total of
$50 \times 2 \times 6 \times 50 = 30{,}000$ repetitions.
GPQA-Diamond yields more eligible tasks than needed in the screening pass, so we
retain $50$ to match the number of eligible HiddenBench tasks, giving a balanced
design with equal task counts per benchmark; the $50$ are drawn uniformly at
random from the eligible GPQA tasks, so that no further selection is made on any
outcome-derived quantity. Because $R = 3$ in the screening pass introduces some
noise, assignments near the inclusion boundary are mildly unstable, which we
accept as a screening cost rather than a final claim.

One consequence of this design should be made explicit for the interpretation of
RQ2. Because the characterisation pass runs only tasks selected by the screening
procedure, raw task difficulty is held roughly constant by construction: it is a
\emph{selection criterion} for the study rather than a factor we vary within it.
The ``difficulty regime'' contrast that RQ2 refers to is therefore operationalised
at the benchmark level, GPQA (a refinement regime, agents mostly start correct)
versus HiddenBench (a correction regime, agents start systematically wrong).
A further interpretive consequence concerns bistability. Because tasks are
retained precisely because their final answer varies across repetitions under the
screening configuration, the prevalence of coexisting attractors reported in
\autoref{sec:results:bistability} is conditioned on this selection and should not
be read as an estimate of motif frequency in an unselected task population; see
\autoref{sec:discussion:limitations} for discussion.

Code, prompts, and run scripts are
available at \url{https://github.com/echtermeyer/LLM-Based-MAS}.
